% 康奈尔听课笔记模板（配色取自王道课件）
\documentclass[10.5pt,a4paper,oneside]{article}
\usepackage[a4paper,top=2.3cm,bottom=2.2cm,left=5.7cm,right=1.5cm,
  marginparwidth=4.25cm,marginparsep=0.55cm,headheight=14pt,footskip=1.1cm]{geometry}
\usepackage{xeCJK}
\usepackage{xeCJKfntef}
\xeCJKDeclareCharClass{FullLeft}{"201C,"2018}
\xeCJKDeclareCharClass{FullRight}{"201D,"2019}
\xeCJKDeclareCharClass{CJK}{"2460->"24FF,"2605,"2606}
\setmainfont{Noto Sans CJK SC}[BoldFont={Noto Sans CJK SC Bold},ItalicFont={Noto Sans CJK SC},ItalicFeatures={FakeSlant=0.18}]
\setsansfont{Noto Sans CJK SC}[BoldFont={Noto Sans CJK SC Bold}]
\setmonofont{DejaVu Sans Mono}[Scale=0.9]
\setCJKmainfont{Noto Sans CJK SC}[BoldFont={Noto Sans CJK SC Bold},ItalicFont={Noto Sans CJK SC},ItalicFeatures={FakeSlant=0.18}]
\setCJKsansfont{Noto Sans CJK SC}[BoldFont={Noto Sans CJK SC Bold}]
\setCJKmonofont{Noto Sans Mono CJK SC}[Scale=0.9]
\linespread{1.32}
\usepackage{amsmath,amssymb}
\usepackage[dvipsnames,table]{xcolor}
\usepackage{graphicx}
\usepackage{adjustbox}
\usepackage{tikz}
\usetikzlibrary{calc}
\usepackage[most]{tcolorbox}
\usepackage{eso-pic}
\usepackage{fancyhdr}
\usepackage{titlesec}
\usepackage{needspace}
\usepackage{enumitem}
\usepackage{xltabular,booktabs,array,colortbl}
\usepackage{marginnote}
\reversemarginpar
\renewcommand*{\raggedleftmarginnote}{}
\renewcommand*{\raggedrightmarginnote}{}
\usepackage{caption}
\usepackage{titletoc}
\usepackage{fancyvrb}
\usepackage{framed}
\usepackage[section]{placeins}
\usepackage{hyperref}
\hypersetup{colorlinks=true,linkcolor=wdblue,urlcolor=wdblue,pdftitle={2.2
课后习题：线性表的顺序表示},pdfauthor={听课笔记}}

% ---------- 课件配色 ----------
\definecolor{wdblue}{HTML}{4472C4}
\definecolor{wdlightblue}{HTML}{5B9BD5}
\definecolor{wdtitle}{HTML}{4A86C8}
\definecolor{wddark}{HTML}{1F3864}
\definecolor{wdorange}{HTML}{ED7D31}
\definecolor{wdgray}{HTML}{A5A5A5}
\definecolor{wdgold}{HTML}{FFC000}
\definecolor{wdgreen}{HTML}{70AD47}
\definecolor{wddeepgreen}{HTML}{3F7F2A}
\definecolor{wdred}{HTML}{C00000}
\definecolor{wdpurple}{HTML}{7030A0}
\definecolor{cuebg}{HTML}{F2F6FC}
\definecolor{codebg}{HTML}{F7F9FC}
\color[HTML]{262626}

% ---------- 行内样式 ----------
\newcommand{\hlred}[1]{\textcolor{wdred}{\bfseries #1}}
\newcommand{\hlblue}[1]{\textcolor{wdblue}{\bfseries #1}}
\newcommand{\hlgreen}[1]{\textcolor{wddeepgreen}{\bfseries #1}}
\newcommand{\hlorange}[1]{\textcolor{wdorange}{\bfseries #1}}
\newcommand{\hlpurple}[1]{\textcolor{wdpurple}{\bfseries #1}}
\newcommand{\uln}[1]{\CJKunderline[format=\color{wdred},thickness=0.9pt,depth=0.28em]{#1}}
\newcommand{\wavy}[1]{\CJKunderwave[format=\color{wdorange}]{#1}}
\newcommand{\hlmark}[1]{\CJKunderline[format=\color{wdgold!55},thickness=0.55em,depth=-0.05em]{#1}}
\newcommand{\tagbox}[2]{\tcbox[on line,colback=#1,colframe=#1,coltext=white,boxsep=0.5pt,left=3pt,right=3pt,top=0.5pt,bottom=0.5pt,arc=2pt,fontupper=\small\bfseries]{#2}}

% ---------- 康奈尔版式背景：左侧线索栏 ----------
\newlength{\cuestripw}\setlength{\cuestripw}{5.15cm}
\AddToShipoutPictureBG{%
  \begin{tikzpicture}[remember picture,overlay]
    \fill[cuebg] (current page.north west) rectangle ($(current page.south west)+(\cuestripw,0)$);
    \draw[wdblue,line width=1.2pt] ($(current page.north west)+(\cuestripw,-1.55cm)$) -- ($(current page.south west)+(\cuestripw,1.55cm)$);
    \node[anchor=west,text=wdblue,font=\small\bfseries] at ($(current page.north west)+(0.75cm,-1.25cm)$) {线索 · 问题};
    \node[anchor=west,text=wdblue,font=\small\bfseries] at ($(current page.north west)+(\cuestripw+0.55cm,-1.25cm)$) {笔记 · 老师主线 + 补充讲解};
  \end{tikzpicture}}

% ---------- 页眉页脚 ----------
\pagestyle{fancy}\fancyhf{}
\renewcommand{\headrulewidth}{0pt}
\fancyhead[R]{\small\color{wdgray}第2章 线性表}
\fancyfoot[C]{\small\color{wdgray}\thepage}
\fancyfoot[R]{\small\color{wdorange}2.2
课后习题：线性表的顺序表示 · 听课笔记}
\fancyfoot[L]{\small\hyperlink{toc}{\textcolor{wdblue}{\ensuremath{\uparrow}\,回到目录}}}
\fancypagestyle{plain}{\fancyhf{}\fancyfoot[C]{\small\color{wdgray}\thepage}\fancyfoot[R]{\small\color{wdorange}2.2
课后习题：线性表的顺序表示 · 听课笔记}}

% ---------- 标题：仿课件的蓝色标题 + 六色条 ----------
\newcommand{\sixbar}{\par\nointerlineskip\vspace{3pt}\centering
  \begin{tikzpicture}\foreach \c [count=\i] in {wdlightblue,wdorange,wdgray,wdgold,wdblue,wdgreen}
     {\fill[\c] ({(\i-1)*1.05},0) rectangle ({(\i-1)*1.05+0.92},0.11);}\end{tikzpicture}\par}
\titleformat{\section}[block]{\centering\Large\bfseries\color{wdtitle}}{}{0pt}{}[\sixbar]
\titlespacing*{\section}{0pt}{16pt}{10pt}
\pretocmd{\section}{\Needspace{8\baselineskip}}{}{}
\titleformat{\subsection}{\Needspace{7\baselineskip}\large\bfseries\color{wddark}}{\textcolor{wdblue}{\rule[-0.15em]{4pt}{1.05em}}\hspace{6pt}}{0pt}{}
\titlespacing*{\subsection}{0pt}{12pt}{5pt}
\titleformat{\subsubsection}{\Needspace{5\baselineskip}\normalsize\bfseries\color{wdblue}}{\textcolor{wdorange}{\large\ensuremath{\blacktriangleright}}\hspace{4pt}}{0pt}{}
\titlespacing*{\subsubsection}{0pt}{8pt}{3pt}
\setcounter{secnumdepth}{0}
\setlength{\parindent}{0pt}\setlength{\parskip}{5pt}
\setlist{nosep,leftmargin=1.5em,itemsep=2pt}
\setlist[itemize,1]{label=\textcolor{wdblue}{\small\ensuremath{\bullet}}}
\setlist[itemize,2]{label=\textcolor{wdorange}{\scriptsize\ensuremath{\blacktriangleright}}}
\setlist[enumerate,1]{label=\textcolor{wdblue}{\bfseries\arabic*.}}
\providecommand{\tightlist}{\setlength{\itemsep}{2pt}\setlength{\parskip}{0pt}}

% ---------- 目录 ----------
\setcounter{tocdepth}{1}
\titlecontents{section}[0pt]{\addvspace{5pt}\bfseries\color{wddark}}{}{}{\titlerule*[0.7pc]{\color{wdgray!60}.}\contentspage}
\titlecontents{subsection}[1.6em]{\addvspace{1pt}\small\color{black!75}}{}{}{\titlerule*[0.7pc]{\color{wdgray!40}.}\contentspage}
\newcommand{\tocback}{\par\vspace{-3pt}{\hfill\footnotesize\hyperlink{toc}{\textcolor{wdgray}{\ensuremath{\uparrow}\,回到目录}}}\par\vspace{-4pt}}

% ---------- 线索栏 ----------
\newcommand{\cueblock}[2]{\Needspace{4\baselineskip}\leavevmode#1{#2}\par\vspace{-\baselineskip}\vspace{-\parskip}}
% 线索直接贴在段首行的左侧（不用 marginnote，避免分页时偶尔落进正文栏）
\newcommand{\cuebox}[1]{\leavevmode\llap{\smash{\begin{minipage}[t]{\marginparwidth}\raggedright\small\linespread{1.18}\selectfont #1\end{minipage}}\hspace{\marginparsep}}}
\newcommand{\cuenote}[1]{\cuebox{{\color{wdblue}\bfseries Q\,}\color{wddark}#1}}
\newcommand{\cuekey}[1]{\cuebox{\color{wdred}\bfseries #1}}

% ---------- 讲解框 ----------
\tcbset{wdbox/.style={enhanced,breakable,arc=2pt,boxrule=0pt,leftrule=3pt,left=7pt,right=7pt,top=5pt,bottom=5pt,
   before skip=7pt,after skip=7pt,fonttitle=\bfseries\small,attach title to upper={\quad},coltitle=#1,
   colframe=#1,colback=#1!6}}
\newtcolorbox{defbox}[1][]{enhanced,breakable,arc=3pt,boxrule=0.8pt,colframe=wdblue,colback=wdblue!4,
   left=8pt,right=8pt,top=6pt,bottom=6pt,fonttitle=\bfseries\small,coltitle=white,colbacktitle=wdblue,
   title={定义#1},before skip=8pt,after skip=8pt}
\newtcolorbox{bridgebox}{wdbox=wdblue,title={承上启下}}
\newtcolorbox{plainbox}{wdbox=wddeepgreen,colframe=wdgreen,colback=wdgreen!8,title={大白话}}
\newtcolorbox{deepbox}{wdbox=wdpurple,title={背后的思想}}
\newtcolorbox{exambox}{wdbox=wdred,title={考点}}
\newtcolorbox{pitfallbox}{wdbox=wdorange,title={易错提醒}}
\newtcolorbox{bookbox}{wdbox=wddark,colback=wdgray!10,colframe=wdgray,title={王道书补充}}
\newtcolorbox{teacherbox}{wdbox=wdorange,colback=wdgold!8,colframe=wdgold,coltitle=wdorange,title={老师原话}}
\newtcolorbox{fixbox}{wdbox=wdred,colback=wdred!4,title={勘误 / 说明}}
% 习题笔记用
\newtcolorbox{questionbox}[1][]{enhanced,breakable,arc=3pt,boxrule=0.8pt,colframe=wddark,colback=white,
   left=8pt,right=8pt,top=6pt,bottom=6pt,fonttitle=\bfseries\small,coltitle=white,colbacktitle=wddark,
   title={#1},before skip=10pt,after skip=6pt}
\newtcolorbox{videobox}{wdbox=wdorange,colback=wdgold!8,colframe=wdgold,coltitle=wdorange,title={视频讲解}}
\newtcolorbox{jiexibox}{wdbox=wddark,colback=wdgray!10,colframe=wdgray,title={书上解析}}
\newtcolorbox{addbox}{wdbox=wdpurple,title={我的补充}}
\newtcolorbox{recallbox}{wdbox=wdblue,title={回顾正课}}
\newtcolorbox{watchbox}{enhanced,breakable,arc=3pt,boxrule=1pt,colframe=wdorange,colback=wdgold!12,
   fonttitle=\bfseries\small,coltitle=white,colbacktitle=wdorange,title={建议去看视频},left=8pt,right=8pt,before skip=8pt,after skip=8pt}
\newtcolorbox{summarybox}{enhanced,arc=3pt,
   boxrule=1pt,colframe=wdblue,colback=white,colbacktitle=wdblue,coltitle=white,fonttitle=\bfseries,
   title={总结区 · 用自己的话复述这一节},left=10pt,right=10pt,top=8pt,bottom=8pt,before skip=14pt}

% ---------- 代码 ----------

\makeatletter
\@ifundefined{Shaded}{}{\renewenvironment{Shaded}{\def\FrameCommand{\fboxsep=5pt\fboxrule=0.5pt\fcolorbox{wdblue!35}{codebg}}%
   \MakeFramed{\advance\hsize-\width\FrameRestore}\small}{\endMakeFramed}}
\makeatother
\fvset{fontsize=\small}

% ---------- 表格 ----------
\renewcommand{\arraystretch}{1.42}
\newcolumntype{Y}{>{\raggedright\arraybackslash}X}

% ---------- 图片 ----------
% 课件截图做成浮动体：LaTeX 会自动用正文填满页面，避免大块空白
\setcounter{topnumber}{3}\setcounter{bottomnumber}{2}\setcounter{totalnumber}{4}
\renewcommand{\topfraction}{0.85}\renewcommand{\bottomfraction}{0.6}\renewcommand{\textfraction}{0.12}\renewcommand{\floatpagefraction}{0.75}
\setlength{\textfloatsep}{10pt plus 3pt minus 2pt}\setlength{\floatsep}{8pt plus 2pt}
\newcommand{\slidefig}[3]{\begin{figure}[!ht]\centering\begin{tcolorbox}[enhanced,width=#2\linewidth,arc=2pt,boxrule=0.6pt,colframe=wdgray!70,
   colback=white,left=2pt,right=2pt,top=2pt,bottom=2pt,drop fuzzy shadow=black!25]\includegraphics[width=\linewidth]{#1}\end{tcolorbox}
   \vspace{-3pt}{\small\color{wdgray}#3}\end{figure}}
\newcommand{\plainfig}[3]{\par\noindent\begin{minipage}{\linewidth}\centering\includegraphics[width=#2\linewidth]{#1}\par\vspace{2pt}{\small\color{wdgray}#3}\end{minipage}\par\vspace{4pt}}

\begin{document}
% ---------- 封面条（仿课件封面） ----------
\thispagestyle{plain}
\begin{tcolorbox}[enhanced,enlarge left by=-4.8cm,width=\linewidth+4.8cm,arc=0pt,boxrule=0pt,
   colback=wddark,interior style={left color=wddark,right color=wdblue},left=18pt,right=18pt,top=14pt,bottom=12pt,
   overlay={\draw[white,line width=2.5pt] ([xshift=10pt,yshift=-10pt]frame.north west) -- ++(0,-0.9cm);
            \draw[white,line width=2.5pt] ([xshift=10pt,yshift=-10pt]frame.north west) -- ++(1.6cm,0);
            \draw[white,line width=2.5pt] ([xshift=-10pt,yshift=10pt]frame.south east) -- ++(0,0.9cm);
            \draw[white,line width=2.5pt] ([xshift=-10pt,yshift=10pt]frame.south east) -- ++(-1.6cm,0);}]
  {\color{white!85!wdlightblue}\small 数据结构 · 第2章
线性表 · 康奈尔听课笔记}\par\vspace{6pt}
  {\color{white}\fontsize{22}{28}\selectfont\bfseries 2.2
课后习题：线性表的顺序表示}\par\vspace{6pt}
  {\color{white!80!wdlightblue}\small 王道书 §2.2.3 试题精选 · §2.2.4
答案与解析（第 18--22 页）· 习题讲解字幕 P7--P9（选择 1--13）}
\end{tcolorbox}
\vspace{4pt}
\begin{tcolorbox}[enhanced,arc=3pt,boxrule=0.8pt,colframe=wdblue,colback=wdblue!3,left=10pt,right=10pt,top=6pt,bottom=6pt,
   colbacktitle=wdblue,coltitle=white,fonttitle=\bfseries\small,title={目录 · 点击标题跳转}]
\hypertarget{toc}{}\makeatletter\@starttoc{toc}\makeatother
\end{tcolorbox}
\hypertarget{ux8fd9ux7ec4ux9898ux5728ux8003ux4ec0ux4e48}{%
\section{这组题在考什么}\label{ux8fd9ux7ec4ux9898ux5728ux8003ux4ec0ux4e48}}

\tocback

\Needspace{4\baselineskip}

\cuenote{13 道题主要围绕哪两件事？}13
道选择题几乎都绕着顺序表的\textbf{一优一缺}转，正课 2.2.1
的“顺序表的特点”和 2.2.2 两节的复杂度就是全部考点：

\begin{itemize}
\tightlist
\item
  \textbf{优点}：\hlred{随机存取}（按序号 \(O(1)\)）、存储密度高；
\item
  \textbf{缺点}：插入、删除要\hlred{移动元素}（平均
  \(O(n)\)）、扩容要整块重新申请。
\end{itemize}

{\rowcolors{2}{wdblue!5}{white}\arrayrulecolor{wdblue!35}
\begin{xltabular}{\linewidth}{>{\hsize=0.474\hsize\raggedright\arraybackslash}X>{\hsize=2.053\hsize\raggedright\arraybackslash}X>{\hsize=0.474\hsize\raggedright\arraybackslash}X}
\rowcolor{wdblue}\textcolor{white}{\bfseries 题号} & \textcolor{white}{\bfseries 考点} & \textcolor{white}{\bfseries 答案} \\
\endfirsthead
01 & 顺序存储的优点：存储密度大 & A \\ \hline
02 & 顺序表 vs 一维数组 & C \\ \hline
03 & 随机存取指什么 & C \\ \hline
04 & 顺序表占用空间与什么无关 & B \\ \hline
05 & 存取第 \(i\) 个及其前驱后继，选存储方式 & D \\ \hline
06 & 按序号存取 + 表尾插删，选存储方式 & A \\ \hline
07 & 数组表示下哪些操作 \(O(1)\) & C \\ \hline
08 & 哪些操作顺序表比链表效率高 & C \\ \hline
09 & 删第 \(i\) 个要移动几个元素 & C \\ \hline
10 & 访问第 \(i\) 个、在第 \(i\) 个插入的复杂度 & C \\ \hline
11 & 哪个运算 \(O(1)\) & C \\ \hline
12 & 扩容申请失败说明什么 & D \\ \hline
13 & 【2023 统考】有序顺序表上平均 \(O(1)\) 的操作 & D \\ \hline
\end{xltabular}\par\vspace{4pt}}

\begin{fixbox}

\textbf{习题讲解视频基本不用看}，内容和书上解析、正课一致，有价值的说法下面都记了。01、13
两段的板书截图已放在对应题下。视频里有几处口误或说得不严的地方，在对应的题里标出。

\end{fixbox}

\hypertarget{ux4e00ux987aux5e8fux8868ux7684ux6982ux5ff5ux4e0eux7279ux70b90104}{%
\section{一、顺序表的概念与特点（01--04）}\label{ux4e00ux987aux5e8fux8868ux7684ux6982ux5ff5ux4e0eux7279ux70b90104}}

\tocback

\Needspace{8\baselineskip}\vspace{22pt}

\cueblock{\cuenote}{\relax{顺序存储比链式存储好在哪？}}

\begin{questionbox}[01]

下列叙述中，（　）是顺序存储结构的优点。

A. 存储密度大　　　　　　B. 插入运算方便\\
C. 删除运算方便　　　　　D. 方便地运用于各种逻辑结构的存储表示

\end{questionbox}

\hlred{\textbf{答案：A}}

\begin{videobox}

老师画内存格子比较：表 \((A,B,C,D)\) 都是 \texttt{char}，各占
1B。顺序存储放在 0～3 号单元，10 个单元能放 10
个字符。链式存储不要求连续，但每个元素后面要跟一个指针域（设
\texttt{int}，4B）指出下一个元素在哪，10 个单元只放得下 2
个字符。同样的空间，顺序存储放得多，\textbf{存储密度大}。

B、C、D
说的都是链式存储的优点：顺序存储插入、删除要移动元素，链式只改指针；链式可以设多个指针域，方便表示各种逻辑结构，顺序存储要求“逻辑相邻
= 物理相邻”，限制了能表示的结构。

\end{videobox}

\slidefig{assets/2.2ex-01.jpg}{0.75}{习题讲解视频板书：同样 10 个单元，顺序存储放 4 个字符还剩 6
格；链式存储每个字符后跟 4B 的 next 下标，只放得下 a、b 两个}

\begin{jiexibox}

顺序表不像链表那样要在结点中存放指针域，因此存储密度大，选项 A
正确。选项 B 和 C 是链表的优点。选项 D
是错误的，比如对于树形结构，顺序表显然不如链表表示起来方便。

\end{jiexibox}

\begin{recallbox}

正课 2.2.1“四、顺序表的特点”：存储密度 = 数据元素本身占的空间 ÷
结点总空间，顺序表为 1，链表小于 1；插入、删除不方便是顺序表的缺点。

\end{recallbox}

\Needspace{8\baselineskip}\vspace{22pt}

\cueblock{\cuenote}{\relax{顺序表和一维数组，在逻辑结构上是一回事吗？}}

\begin{questionbox}[02]

下列关于顺序表的叙述中，正确的是（　）。

A.
顺序表可以利用一维数组表示，因此顺序表与一维数组在逻辑结构上是相同的\\
B. 在顺序表中，逻辑上相邻的元素物理位置上不一定相邻\\
C. 顺序表和一维数组一样，都可以进行随机存取\\
D. 在顺序表中，每个元素的类型不必相同

\end{questionbox}

\hlred{\textbf{答案：C}}

\begin{videobox}

先建立认知：顺序表\textbf{逻辑上}是线性表，\textbf{存储上}是顺序存储。

\begin{itemize}
\tightlist
\item
  A
  前半句对，后半句错：一维数组本质是一种存储方式，用它存的不一定是线性表，后面学的二叉树也能用数组存。
\item
  B 错：顺序存储要求连续存放，逻辑上相邻的元素物理上\textbf{必然}相邻。
\item
  C 对：随机存取由存储结构决定，数组按 \(a_0\)
  的地址加偏移量直接算出任意元素的地址，顺序表就是用一维数组存的，当然也能随机存取。
\item
  D 错：线性表要求元素类型相同，顺序表是线性表，元素类型必须相同。
\end{itemize}

\end{videobox}

\begin{jiexibox}

顺序表是顺序存储的线性表，表中所有元素的类型必须相同，且必须连续存放。一维数组中的元素可以不连续存放；此外，栈、队列和树等逻辑结构也可利用一维数组表示，但它与顺序表不属于相同的逻辑结构。在顺序表中，逻辑上相邻的元素物理位置上也相邻。

\end{jiexibox}

\begin{addbox}

A 的错因和 2.1 课后习题 02 的
D（邻接表）是同一类：\textbf{逻辑结构和存储结构不在一个层面}。顺序表 =
线性表（逻辑）+ 顺序存储（存储）；一维数组只说明了存储方式。

\end{addbox}

\Needspace{8\baselineskip}\vspace{22pt}

\cueblock{\cuenote}{\relax{“随机存取”到底指什么？}}

\begin{questionbox}[03]

通常说顺序表具有随机存取的特性，指的是（　）。

A. 查找值为 \(x\) 的元素的时间与顺序表中元素个数 \(n\) 无关\\
B. 查找值为 \(x\) 的元素的时间与顺序表中元素个数 \(n\) 有关\\
C. 查找序号为 \(i\) 的元素的时间与顺序表中元素个数 \(n\) 无关\\
D. 查找序号为 \(i\) 的元素的时间与顺序表中元素个数 \(n\) 有关

\end{questionbox}

\hlred{\textbf{答案：C}}

\begin{videobox}

随机存取两层意思：① 访问任何位置的时间都一样，访问 \(a_0\) 和
\(a_{100}\) 一样快；②
不用从头遍历，直接访问。原因是顺序表开辟的是一块连续空间，每个元素占 4B
的话，\(a_i\) 的地址就是首地址 \(+4i\)，给出序号马上算出地址。

A 错：按值查找和 \(n\) 有关，即使表有序、用以后学的折半查找也一样。B
本身说得对，但它说的不是随机存取，题目问的是随机存取指什么。

\end{videobox}

\begin{jiexibox}

随机存取是指在 \(O(1)\) 的时间访问下标为 \(i\)
的元素，所需时间与顺序表中的元素个数 \(n\) 无关。

\end{jiexibox}

\begin{recallbox}

习题讲解老师的“首地址 \(+4i\)”里 \(i\) 是\textbf{下标}（从 0
开始）；正课 2.2.1 的公式
\(LOC(a_i)=LOC(L)+(i-1)\times\texttt{sizeof(ElemType)}\) 里 \(i\)
是\textbf{位序}。两者是一回事，算地址时先看清题目给的是哪个。按位查找、按值查找的
\(O(1)\) 和 \(O(n)\) 见正课 2.2.2\_2。

\end{recallbox}

\Needspace{8\baselineskip}\vspace{22pt}

\cueblock{\cuenote}{\relax{顺序表占多大空间，由哪些因素决定？}}

\begin{questionbox}[04]

一个顺序表所占用的存储空间大小与（　）无关。

A. 表的长度　　　　　　　B. 元素的存放顺序\\
C. 元素的类型　　　　　　D. 元素中各字段的类型

\end{questionbox}

\hlred{\textbf{答案：B}}

\begin{videobox}

看结构体定义：\texttt{length} 只是一个
\texttt{int}，大小主要由数组决定，数组多大取决于元素个数和每个元素多大。A
表越长占得越多；C 元素类型不同，每个元素大小不同；D
元素是结构体时，各字段类型决定单个元素的大小，和 C 是一个意思。B
存放顺序变了，总空间不变。

\end{videobox}

\begin{jiexibox}

顺序表所占的存储空间 = 表长 ×
sizeof（元素的类型），表长和元素的类型显然会影响存储空间的大小。若元素为结构体类型，则元素中各字段的类型也会影响存储空间的大小。

\end{jiexibox}

\begin{fixbox}

习题讲解视频举 C 的例子时说“100 个字符型变量\ldots\ldots 占 400
个字节”，是口误：100 个 \texttt{char} 占 100B，100 个
\texttt{int}（4B）才占 400B。

\end{fixbox}

\hypertarget{ux4e8cux6309ux9700ux6c42ux9009ux5b58ux50a8ux65b9ux5f0f0506}{%
\section{二、按需求选存储方式（05--06）}\label{ux4e8cux6309ux9700ux6c42ux9009ux5b58ux50a8ux65b9ux5f0f0506}}

\tocback

\Needspace{8\baselineskip}\vspace{22pt}

\cueblock{\cuenote}{\relax{常用操作是按序号存取，选哪种存储？}}

\begin{questionbox}[05]

若线性表最常用的操作是存取第 \(i\)
个元素及其前驱和后继元素的值，为了提高效率，应采用（　）的存储方式。

A. 单链表　　　　B. 双链表　　　　C. 循环单链表　　　　D. 顺序表

\end{questionbox}

\hlred{\textbf{答案：D}}

\begin{videobox}

把题干翻译一下：存取
\(a_i\)、\(a_{i-1}\)、\(a_{i+1}\)，都是\textbf{按序号}定位，就是随机存取，只有顺序存储能
\(O(1)\)
做到。不管单链表、双链表还是循环单链表，只要是链式存储，找序号为 \(i\)
的元素都得从头遍历，时间取决于它在哪。

\end{videobox}

\begin{jiexibox}

题干实际要求能最快存取第 \(i-1\)、\(i\) 和 \(i+1\) 个元素值。选项
A、B、C 都只能从头结点依次顺序查找，时间复杂度为
\(O(n)\)；只有顺序表可以按序号随机存取，时间复杂度为 \(O(1)\)。

\end{jiexibox}

\Needspace{8\baselineskip}\vspace{22pt}

\cueblock{\cuenote}{\relax{“在最后进行插入、删除”会不会拖累顺序表？}}

\begin{questionbox}[06]

一个线性表最常用的操作是存取任意一个指定序号的元素并在最后进行插入、删除操作，则利用（　）存储方式可以节省时间。

A. 顺序表　　　　　　　　　B. 双链表\\
C. 带头结点的循环双链表　　D. 循环单链表

\end{questionbox}

\hlred{\textbf{答案：A}}

\begin{videobox}

第一个条件“存取指定序号的元素”就已经锁定顺序表。有同学会犹豫：顺序表的缺点不是插入删除要移动大量元素吗？但那说的是\textbf{平均情况}。在表尾插入、删除正好是最好情况，一个元素都不用移动。

\end{videobox}

\begin{jiexibox}

只有顺序表可以按序号随机存取，且在最后进行插入和删除操作时不需要移动任何元素。

\end{jiexibox}

\begin{recallbox}

正课
2.2.2\_1：插入最好情况是插在表尾（\(i=n+1\)），删除最好情况是删表尾（\(i=n\)），都是
\(O(1)\)。“要移动大量元素”指的是平均 \(O(n)\)。

\end{recallbox}

\hypertarget{ux4e09ux5404ux64cdux4f5cux7684ux65f6ux95f4ux590dux6742ux5ea60711}{%
\section{三、各操作的时间复杂度（07--11）}\label{ux4e09ux5404ux64cdux4f5cux7684ux65f6ux95f4ux590dux6742ux5ea60711}}

\tocback

\Needspace{8\baselineskip}\vspace{22pt}

\cueblock{\cuenote}{\relax{数组表示的线性表，哪些操作是 \(O(1)\)？}}

\begin{questionbox}[07]

在 \(n\) 个元素的线性表的数组表示中，时间复杂度为 \(O(1)\)
的操作是（　）。

Ⅰ. 访问第 \(i\,(1\le i\le n)\) 个结点和求第 \(i\,(2\le i\le n)\)
个结点的直接前驱\\
Ⅱ. 在最后一个结点后插入一个新的结点\\
Ⅲ. 删除第 1 个结点\\
Ⅳ. 在第 \(i\,(1\le i\le n)\) 个结点后插入一个结点\\
A. Ⅰ　　　　B. Ⅱ、Ⅲ　　　　C. Ⅰ、Ⅱ　　　　D. Ⅰ、Ⅱ、Ⅲ

\end{questionbox}

\hlred{\textbf{答案：C}}

\begin{videobox}

“数组表示”就是顺序表。

\begin{itemize}
\tightlist
\item
  Ⅰ：给序号 \(i\) 直接定位，前驱就是序号
  \(i-1\)，也是随机访问，\(O(1)\)。
\item
  Ⅱ：表尾后插入，算出地址直接放，不用移动，\(O(1)\)。
\item
  Ⅲ：删第 1 个，后面 \(n-1\) 个都要前移，\(O(n)\)。
\item
  Ⅳ：插在任意位置，平均要移动约 \(\dfrac n2\)
  个，\(O(n)\)。老师说平均到底是 \(\dfrac n2\)
  还是别的不重要，不影响数量级。
\end{itemize}

\end{videobox}

\begin{jiexibox}

对说法Ⅰ，解析略；对说法Ⅱ，在最后位置插入新结点不需要移动元素，时间复杂度为
\(O(1)\)；对说法Ⅲ，被删结点后的结点需要依次前移，时间复杂度为
\(O(n)\)；对说法Ⅳ，需要后移 \(n-i\) 个结点，时间复杂度为 \(O(n)\)。

\end{jiexibox}

\begin{addbox}

Ⅳ 是“在第 \(i\) 个\textbf{之后}”插入，即插到第 \(i+1\) 个位置，要移动第
\(i+1\)～\(n\) 个，共 \(n-i\) 个（用正课的
\(n-i'+1\)，\(i'=i+1\)，结果一样）。Ⅱ 就是 Ⅳ 在 \(i=n\) 时的特例。

\end{addbox}

\Needspace{8\baselineskip}\vspace{22pt}

\cueblock{\cuenote}{\relax{两种存储复杂度一样时，还能比效率吗？}}

\begin{questionbox}[08]

设线性表有 \(n\)
个元素，严格来说，以下操作中，（　）在顺序表上实现要比在链表上实现的效率高。

Ⅰ. 输出第 \(i\,(1\le i\le n)\) 个元素值\\
Ⅱ. 交换第 3 个元素与第 4 个元素的值\\
Ⅲ. 顺序输出这 \(n\) 个元素的值\\
A. Ⅰ　　　　B. Ⅰ、Ⅲ　　　　C. Ⅰ、Ⅱ　　　　D. Ⅱ、Ⅲ

\end{questionbox}

\hlred{\textbf{答案：C}}

\begin{videobox}

\begin{itemize}
\tightlist
\item
  Ⅰ：顺序表 \(O(1)\) 直接输出；链表要从头遍历到第 \(i\) 个，\(O(n)\)。
\item
  Ⅱ：顺序表用一个临时变量，三句赋值就换完。链表要先从头数到第 3
  个，再改指针。注意 \uln{两者其实都是 \(O(1)\)}：位置已经定死是第 3、4
  个，和表长无关。但顺序表执行的语句更少，效率仍然更高。老师的话：执行 1
  条语句和执行 100 条语句都是 \(O(1)\)，效率却不一样。
\item
  Ⅲ：两者都要遍历整张表，开销近似，不选。
\end{itemize}

\end{videobox}

\begin{jiexibox}

对说法Ⅱ，顺序表只需 3 次交换操作；链表需要分别找到两个结点前驱，第 4
个结点断链后再插入到第 2
个结点后，效率较低。对说法Ⅲ，需要依次顺序访问每个元素，时间复杂度相同。

\end{jiexibox}

\begin{addbox}

这题的“严格来说”就是在提示：\textbf{时间复杂度相同 ≠
效率相同}。复杂度只看数量级，常数大小也算效率。

\end{addbox}

\Needspace{8\baselineskip}\vspace{22pt}

\cueblock{\cuenote}{\relax{删第 \(i\) 个元素，要移动几个？}}

\begin{questionbox}[09]

在一个长度为 \(n\) 的顺序表中删除第 \(i\,(1\le i\le n)\)
个元素时，需向前移动（　）个元素。

A. \(n\)　　　　B. \(i-1\)　　　　C. \(n-i\)　　　　D. \(n-i+1\)

\end{questionbox}

\hlred{\textbf{答案：C}}

\begin{videobox}

删掉 \(a_i\) 后，逻辑上 \(a_{i-1}\) 的后继变成 \(a_{i+1}\)，存储上就要把
\(a_i\) 后面的元素都前移一位。要移的是 \(a_{i+1}\)～\(a_n\)，个数
\(n-(i+1)+1=n-i\)。

\end{videobox}

\begin{jiexibox}

需要将元素 \(a_{i+1}\sim a_n\) 依次前移一位，共移动 \(n-(i+1)+1=n-i\)
个元素。

\end{jiexibox}

\begin{recallbox}

数个数的通用方法：从 \(a_p\) 到 \(a_q\) 共 \(q-p+1\) 个。插入第 \(i\)
个位置要后移 \(a_i\sim a_n\)，共 \(n-i+1\) 个（D 就是这个干扰项）。正课
2.2.2\_1 的两张循环次数表：删除 \(i=1\) 循环 \(n-1\) 次、插入 \(i=1\)
循环 \(n\) 次，代进去正好对得上。

\end{recallbox}

\Needspace{8\baselineskip}\vspace{22pt}

\cueblock{\cuenote}{\relax{访问第 \(i\) 个和在第 \(i\) 个插入，复杂度各是多少？}}

\begin{questionbox}[10]

对于顺序表，访问第 \(i\) 个位置的元素和在第 \(i\)
个位置插入一个元素的时间复杂度为（　）。

A. \(O(n),\,O(n)\)　　　B. \(O(n),\,O(1)\)　　　C.
\(O(1),\,O(n)\)　　　D. \(O(1),\,O(1)\)

\end{questionbox}

\hlred{\textbf{答案：C}}

\begin{videobox}

访问：随机存取，\(O(1)\)。插入：\(n\) 个元素有几个插入位置？相当于 \(n\)
个点把一条直线分成几段，答案 \(n+1\) 段，即 \(n+1\) 个位置。插在 \(b_1\)
前移 \(n\) 个，插在 \(b_1,b_2\) 之间移 \(n-1\) 个\ldots\ldots 插在
\(b_n\) 后移 0 个，平均
\(\dfrac{0+1+\cdots+n}{n+1}=\dfrac n2\)，\(O(n)\)。

\end{videobox}

\begin{jiexibox}

在第 \(i\) 个位置插入一个元素，需要移动 \(n-i+1\) 个元素，时间复杂度为
\(O(n)\)。

\end{jiexibox}

\Needspace{8\baselineskip}\vspace{22pt}

\cueblock{\cuenote}{\relax{排序、插入、删除、修改，哪个是 \(O(1)\)？}}

\begin{questionbox}[11]

对于顺序存储的线性表，其算法时间复杂度为 \(O(1)\) 的运算应该是（　）。

A. 将 \(n\) 个元素从小到大排序\\
B. 删除第 \(i\,(1\le i\le n)\) 个元素\\
C. 改变第 \(i\,(1\le i\le n)\) 个元素的值\\
D. 在第 \(i\,(1\le i\le n)\) 个元素后插入一个新元素

\end{questionbox}

\hlred{\textbf{答案：C}}

\begin{videobox}

C：给了序号 \(i\)，直接给 \(a_i\) 赋值，\(O(1)\)。B、D：删除平均移
\(\dfrac{n-1}{2}\) 个、插入平均移 \(\dfrac n2\) 个，都是
\(O(n)\)。A：排序第 8 章才学，老师说最好的排序算法也要 \(n\log n\)
量级。

\end{videobox}

\begin{jiexibox}

对 \(n\) 个元素进行排序的时间复杂度最小也要
\(O(n)\)（初始有序时），通常为 \(O(n\log_2 n)\) 或 \(O(n^2)\)，通过第 8
章学习后会更容易理解。选项 B 和 D
显然错误。顺序表支持按序号的随机存取方式。

\end{jiexibox}

\begin{fixbox}

习题讲解老师说排序“最好也要 \(n\log n\)“，书上解析说”最小也要
\(O(n)\)（初始有序时）”，两种说法角度不同：书说的是初始已经有序这种最好情况；\(n\log n\)
是常见高效排序的一般水平。这里只需要知道排序不可能 \(O(1)\)，细节第 8
章再说。

\end{fixbox}

\begin{addbox}

“改”就是先“查”再赋值（2.1
正课老师说的“改的第一步就是查”）。顺序表按位查是
\(O(1)\)，所以按位改也是 \(O(1)\)。

\end{addbox}

\hypertarget{ux56dbux6269ux5bb912}{%
\section{四、扩容（12）}\label{ux56dbux6269ux5bb912}}

\tocback

\Needspace{8\baselineskip}\vspace{22pt}

\cueblock{\cuenote}{\relax{扩容失败，是缺 \(m\) 个空间，还是缺 \(n+m\) 个连续空间？}}

\begin{questionbox}[12]

顺序表的插入算法中，当 \(n\) 个空间已满时，可再申请增加分配 \(m\)
个空间，若申请失败，则说明系统没有（　）可分配的存储空间。

A. \(m\) 个　　　　　　　　B. \(m\) 个连续\\
C. \(n+m\) 个　　　　　　　D. \(n+m\) 个连续

\end{questionbox}

\hlred{\textbf{答案：D}}

\begin{videobox}

两个关键点：① 顺序表基于数组，申请的空间必须\textbf{连续}，排除 A、C；②
扩容时申请的是一整块\textbf{新的} \(n+m\) 大小的空间，而不是在原来 \(n\)
个后面接 \(m\) 个。原来那块不能复用。不清楚第 ② 点就会错选 B。

\end{videobox}

\begin{jiexibox}

顺序存储需要连续的存储空间，在申请时需申请 \(n+m\)
个连续的存储空间，然后将线性表原来的 \(n\) 个元素复制到新申请的 \(n+m\)
个连续的存储空间的前 \(n\) 个单元。

\end{jiexibox}

\begin{recallbox}

这就是正课 2.2.1 的 \texttt{IncreaseSize}：先用 \texttt{p}
记住旧空间，\texttt{malloc} 一块 \texttt{MaxSize+len}
的\textbf{新}空间，把 \(n\) 个元素复制过去，再
\texttt{free(p)}。所以申请失败说明找不到 \(n+m\) 个连续单元。

\end{recallbox}

\begin{fixbox}

习题讲解视频说“无论是动态分配还是静态分配，申请新空间时都是全新的内存单元”，说得不严：按正课
2.2.1，\uln{静态分配的数组长度是固定的，满了就没办法扩}，能扩容的只有动态分配。

\end{fixbox}

\hypertarget{ux4e942023-ux7edfux8003ux771fux989813}{%
\section{五、2023
统考真题（13）}\label{ux4e942023-ux7edfux8003ux771fux989813}}

\tocback

\Needspace{8\baselineskip}\vspace{22pt}

\cueblock{\cuenote}{\relax{有序顺序表上，查找、插入、删除、按位获取，谁是平均 \(O(1)\)？}}

\begin{questionbox}[13]

【2023 统考真题】在下列对顺序存储的有序表（长度为
\(n\)）实现给定操作的算法中，平均时间复杂度为 \(O(1)\) 的是（　）。

A. 查找包含指定值元素的算法　　　　B. 插入包含指定值元素的算法\\
C. 删除第 \(i\,(1\le i\le n)\) 个元素的算法　　D. 获取第
\(i\,(1\le i\le n)\) 个元素的算法

\end{questionbox}

\hlred{\textbf{答案：D}}

\begin{videobox}

划重点：“顺序存储”“有序表”，考试中可以把顺序表就当成数组。

\begin{itemize}
\tightlist
\item
  A：可以折半查找（条件：有序 + 顺序存储，都满足），也要 \(O(\log n)\)。
\item
  B：例如在 1、2、3、4\ldots 中插入
  3.5，先查位置（折半，\(O(\log n)\)），再把后面的元素都后移一格，最坏插在最前面，要移
  \(n\) 个，合起来 \(O(n)\)。
\item
  C：删除后面的都前移，\(O(n)\)。
\item
  D：“获取”就是读，\texttt{x\ =\ a{[}i{]}} 一条语句，\(O(1)\)。
\end{itemize}

\end{videobox}

\slidefig{assets/2.2ex-13.jpg}{0.75}{习题讲解视频板书：有序表的定义（数据元素按元素值大小顺序排列的线性表）和四个选项的分析}

\begin{jiexibox}

对于顺序存储的有序表，查找指定值元素可以采用顺序查找法或折半查找法，平均时间复杂度最少为
\(O(\log_2 n)\)。插入指定值元素需要先找到插入位置，然后将该位置及之后的元素依次后移一个位置，最后将指定值元素插入到该位置，平均时间复杂度为
\(O(n)\)。删除第 \(i\)
个元素需要将该元素之后的全部元素依次前移一个位置，平均时间复杂度为
\(O(n)\)。获取第 \(i\)
个元素只需直接根据下标读取对应的数组元素即可，时间复杂度为 \(O(1)\)。

\end{jiexibox}

\begin{addbox}

\textbf{习题讲解老师补充的删除技巧}：如果表是\textbf{无序}的，且允许改变元素位置，删除某个元素可以把它和最后一个元素\textbf{交换}，再
\texttt{length-\/-}，后面没有元素要前移，\(O(1)\)。本题是有序表，交换会破坏有序，所以不能用，C
仍是 \(O(n)\)。

另外，题目问的是\textbf{平均}复杂度，习题讲解视频分析 B、C
时用的是最坏情况；本题平均和最坏都是 \(O(n)\)，结论不变。

\end{addbox}

\begin{recallbox}

正课 2.2.2\_2 老师说过：元素有序时有更高效的查找算法（第 7
章），这道题的 A 就是折半查找。正课 2.2.2
的王道书考点追踪“顺序表上操作的时间复杂度分析（2023）”说的就是这题。

\end{recallbox}

\hypertarget{ux516dux603bux7ed3}{%
\section{六、总结}\label{ux516dux603bux7ed3}}

\tocback

\begin{summarybox}

\textbf{这 13 题就考两件事。}

\textbf{随机存取是顺序表最大的优点}：按序号（位序 /
下标）访问、修改、求前驱后继都是 \(O(1)\)，和 \(n\)
无关；按值查找不是随机存取，和 \(n\)
有关。链式存储（单链表、双链表、循环链表）一律要从头找，做不到。顺序表的另一个优点是存储密度大（没有指针域）。

\textbf{插入删除要移动元素是最大的缺点}：插第 \(i\) 个移 \(n-i+1\)
个，删第 \(i\) 个移 \(n-i\) 个，平均都是
\(O(n)\)；但表尾插删一个都不用移，是 \(O(1)\)。扩容要重新申请 \(n+m\)
个\textbf{连续}的新空间再复制。

两个容易忽略的点：复杂度相同不代表效率相同（08）；有序表能折半查找，但插入仍是
\(O(n)\)（13）。

\textbf{自测}：① 长度为 \(n\) 的顺序表，在第 \(i\)
个元素\textbf{之后}插入要移动几个？②
为什么“顺序表与一维数组在逻辑结构上相同”是错的？③ 动态分配的顺序表扩容
\(m\) 个，要申请多少个连续空间？④ 无序顺序表允许打乱顺序时，怎样
\(O(1)\) 删除一个元素？

\end{summarybox}
\end{document}
